% SPDX-License-Identifier: CC-BY-4.0
% P5 autonomous fault observability & Byzantine integrity manuscript.
% Target: NAVIGATION (ION) / IEEE TAES; conference flag ION GNSS+ 2026.
% Build: tectonic main.tex   (or pdflatex + bibtex + pdflatex x2)
\documentclass[11pt]{article}

\usepackage[utf8]{inputenc}
\usepackage[T1]{fontenc}
\usepackage{lmodern}
\usepackage[margin=1in]{geometry}
\setlength{\emergencystretch}{3em}
\usepackage{amsmath}
\usepackage{amssymb}
\usepackage{amsthm}
\usepackage{bm}
\usepackage{graphicx}
\usepackage{booktabs}
\usepackage{array}
\usepackage{siunitx}
\usepackage{xcolor}
\usepackage[colorlinks=true,linkcolor=black,citecolor=black,urlcolor=blue]{hyperref}
\usepackage[numbers,sort&compress]{natbib}
\usepackage{orcidlink}

\graphicspath{{./}{figures/}}
\setlength{\abovecaptionskip}{9pt}
\setlength{\belowcaptionskip}{6pt}

% Theorem environments
\newtheorem{theorem}{Theorem}
\newtheorem{lemma}{Lemma}
\newtheorem{proposition}{Proposition}
\newtheorem{corollary}{Corollary}
\theoremstyle{definition}
\newtheorem{definition}{Definition}
\newtheorem{remark}{Remark}

% Notation macros
\newcommand{\code}[1]{\texttt{#1}}
\newcommand{\R}{\mathbb{R}}
\newcommand{\Gmat}{\bm{G}}                        % geometry (design) matrix
\newcommand{\Wmat}{\bm{W}}                        % weight matrix Sigma^{-1}
\newcommand{\Pperp}{\bm{P}^{\perp}}               % parity projector
\newcommand{\bvec}{\bm{b}}                        % fault (bias) vector
\newcommand{\xvec}{\bm{x}}                        % state vector
\newcommand{\thetav}{\bm{\theta}}                 % datum vector
\newcommand{\Bmat}{\bm{B}}                        % peer fault signature block
\newcommand{\Gbar}{\bar{\bm{G}}}                  % effective (parity-projected) signature
\newcommand{\Nsp}{\mathcal{N}}                    % null space
\newcommand{\dtau}{\delta\tau}                    % timescale offset gauge
\newcommand{\dalpha}{\delta\alpha}               % timescale rate gauge
\newcommand{\rng}{\operatorname{range}}
\newcommand{\Sig}{\bm{\Sigma}}                    % covariance
\newcommand{\spark}{\operatorname{spark}}
\DeclareMathOperator{\rank}{rank}
\DeclareMathOperator{\MDB}{MDB}

% Honesty tags: Cited (prior art), Proven (our algebra), Validated (external oracle), Modelled (representative)
\newcommand{\Cited}{\textsc{Cited}}
\newcommand{\Proven}{\textsc{Proven}}
\newcommand{\Modelled}{\textsc{Modelled}}
\newcommand{\Validated}{\textsc{Validated}}

\title{\textbf{Autonomous fault observability and Byzantine integrity of a
sparse multi-provider lunar constellation: the datum--timescale gauge as the
undetectable-fault subspace, and a geometric secure-identification bound}}

\author{Chakshu Baweja\,\orcidlink{0009-0008-2098-0751}\\
\small Ashforde O\"U, Tallinn, Estonia \\
\small \texttt{contact@ashforde.org}}

\date{\today}

\begin{document}
\maketitle

\begin{abstract}
An autonomous lunar positioning, navigation, and timing (PNT) constellation, operating in
an Earth-link-loss contingency, must certify its own integrity from \emph{internal}
inter-node measurements alone. We ask, for a sparse multi-provider network, exactly which
of its own faults it can see and how many untrusted (Byzantine) peers it tolerates, and we
separate throughout what we cite, prove, validate, and model. Our organizing result is a
unification: the \emph{single} rank deficiency of the self-referential network Jacobian
$\Gmat$ produces both faces of the problem at once---the datum$\oplus$timescale
\emph{gauge} $\Nsp(\Gmat)\subset\R^{40}$ (the absolute frame and the common clock offset
and rate are unobservable without an external tie) and the \emph{blind fault subspace}
$\rng(\Gmat)\subset\R^{84}$ (the biases a residual test cannot see)---so an undetectable
fault is precisely one absorbed into a state error defined only \emph{modulo} the gauge. On
this footing we prove three things. (i)~A fault is undetectable if and only if it lies in
$\rng(\Gmat)$; the classical parity characterization (\Cited{}) survives the rank-deficient,
gauge-carrying $\Gmat$ through the pseudo-inverse projector (\Proven{}). (ii)~A geometric
secure-identification bound: a peer coalition of size $\le f$ is detectable when the
parity-projected peer blocks are full column rank (block-spark $>f$) and identifiable when
block-spark $>2f$---the coding / sparse-observability bound (\Cited{}) instantiated on the
coupled-datum lunar geometry, \emph{not} a new bound, and explicitly distinct from the
message-passing consensus threshold. (iii)~An autonomous-holdover floor: the common clock
rate is an irreducible temporal-gauge free parameter that no internal monitor removes, and
no finite alert limit protects the $\rng(\Gmat)$ class without an external tie. Exactly one
claim is externally \Validated{}: the fault-observability linear-algebra pipeline---parity
projector, detectability norm, minimum detectable bias $\MDB=\sqrt{\lambda_0/(\bm c^\top\Wmat\Pperp\bm c)}$,
and block-spark rank---reproduces an independent numpy/scipy oracle to below $10^{-3}$ on
real DE440-orientation rows (defect $8$, $\rank\Gmat=32$, parity dimension $52$). The
geometry, the tolerance counts read as operational guarantees, and the holdover magnitude
are \Modelled{}; the near prior art---rigidity-based fault \emph{detection}, secure state
estimation, and LiAISON crosslink observability---is \Cited{} and its boundary stated. Four
design laws follow: engineer redundancy to raise the Byzantine tolerance, spread each node's
measurements so no peer owns its own block, supply an external time/frame tie for the
common-rate and common-provider classes no internal monitor can see, and accept the residual
protection gap.

\end{abstract}

\noindent\textbf{Keywords:} lunar PNT; autonomous integrity monitoring; fault
detection and exclusion; RAIM; parity space; protection level; Byzantine fault
tolerance; secure state estimation; sparse observability; datum defect; gauge
freedom; timescale holdover; LunaNet; Moonlight; inter-satellite ranging; DE440

\vspace{0.5em}
\section{Introduction}
\label{sec:intro}

The next decade will place several independent positioning, navigation, and timing (PNT)
services around the Moon---NASA's LunaNet and its Lunar Communications Relay and Navigation
Systems, ESA's Moonlight programme and its Lunar Communications and Navigation Service, and
their national counterparts \citep{israel2020lunanet,fienga2024moonlight,lnis2023}. Concretely,
NASA's Lunar Augmented Navigation Service (LANS) is disseminated by LunaNet Service Providers
(LNSPs) over an Augmented Forward Signal (AFS), with ESA's Moonlight providing the analogous
European capability. These
services are nominally \emph{Earth-referenced}: their frames and timescales are tied, at
some cadence, to terrestrial reference systems and to Terrestrial Time through Earth--Moon
links. The design driver of the present paper is the \emph{contingency} mode. When the
Earth link is lost or degraded---the case in which integrity and timekeeping matter
most---such a constellation must run \emph{autonomously}: it must certify its own integrity,
detect and exclude faulty members, and hold its timescale, using only the measurements its
nodes make of \emph{each other}. We therefore ask a sharp question about a self-referential
lunar network: \emph{what can it see about its own faults, and how many untrusted peers can
it tolerate, as a function of geometry?}

This paper answers that question and, in doing so, closes a companion arc. P1 classifies the
\emph{spatial} datum defect of a lunar network \citep{baweja2026datum}; P2 quantifies the
irreducible cross-provider interoperability floor \citep{baweja2026interop}; P3 identifies
the common-mode error class an integrity monitor cannot see \citep{baweja2026commonmode};
and P4 adds \emph{time} and the coupled frame$\oplus$timescale gauge \citep{baweja2026gauge}.
The present paper (P5) turns that gauge machinery on \emph{fault observability} and shows
that these gauges and the fault-hiding directions of an autonomous integrity monitor are two
faces of a single rank deficiency of the network Jacobian~$\Gmat$: the spatial datum (P1) and
coupled timescale (P4) gauges are the parameter-space null space $\Nsp(\Gmat)$, whereas the
P3 common-mode class is a subset of the measurement-space range $\rng(\Gmat)$---linked by that
one deficiency, not equal to one another. We are deliberate throughout about the status of every
statement---cited, proven, validated, or modelled---because the honest content of the paper
is a small number of proven structural results and exactly one externally validated pipeline,
set against a large body of correctly attributed prior art.

\paragraph{The prior art, and where our contribution begins.}
Autonomous integrity monitoring is not a new subject, and we bound our novelty up front. The
parity-space characterization of undetectable faults, the minimum detectable bias (MDB), and
the protection level (PL) are classical RAIM/ARAIM \citep{sturza1988,brown1992,parkinson1988,%
baarda1968,teunissen2006,blanch2015,joerger2014}; we cite them and claim none. Cooperative and
collaborative integrity monitoring---solution-separation fault detection and exclusion
\citep{joerger2014} carried across a network of receivers or vehicles \citep{bhamidipati2019coop}---is
equally prior art, but it is empirical and Earth-anchored and does not characterize the
frame$\oplus$timescale gauge as a fault-hiding subspace, which is what we add. The
sparse-observability / error-correcting-code bound that governs how many arbitrary sensor
faults an estimator can survive is the secure-state-estimation and compressed-sensing
literature \citep{fawzi2014,shoukry2016,candes2005,donoho2003}; the bound is theirs. The
Byzantine (arbitrary, colluding) fault \emph{model} is Lamport--Shostak--Pease
\citep{lamport1982}, cited for the model only---emphatically \emph{not} for its $3f+1$
message-passing consensus threshold, which is a different problem class we distinguish
carefully in Section~\ref{sec:byzantine}. Closest of all is the recent, excellent work on
\emph{autonomous constellation fault monitoring via inter-satellite links using rigidity}
\citep{iiyama2024} and its clock-fault follow-on using Euclidean distance matrices
\citep{iiyama2025edm}, built on the rigidity theory of network localization
\citep{aspnes2006}; and the LiAISON result that absolute state is recoverable from crosslinks
given a dynamically-unique orbit \citep{hillborn2007}. This prior art occupies most of the
territory of ``autonomous lunar ISL fault detection.'' Our contribution is the complement of
it, and we state the boundary plainly:
\begin{itemize}
  \item \textbf{Versus rigidity-based monitoring \citep{iiyama2024,iiyama2025edm}.} The
  rigidity approach and its Euclidean-distance-matrix clock-fault follow-on are a
  simulation-validated \emph{algorithm}---a singular-value / GCEDM test on the inter-node
  geometry that detects faults and identifies a \emph{single} faulty satellite, clock jumps
  included. Our contribution is the complement, in three respects. First, we give an
  \emph{algebraic characterization}---undetectable $\iff\rng(\Gmat)$, with provable
  minimum-detectable-bias and block-spark bounds---of \emph{which} faults are provably
  invisible, rather than a validated detector. Second, we bound worst-case
  \emph{colluding-coalition (Byzantine)} identification (block-spark $>2f$,
  Section~\ref{sec:byzantine}), a stronger adversary than a single culprit. Third, we
  characterize the common clock offset and rate as an \emph{unobservable gauge with a holdover
  floor} (Section~\ref{sec:holdover}), together with the coupled datum$\oplus$timescale
  unification and the provider common/differential split---none of which the rigidity papers
  address, even where they treat clock jumps. We concede the deep connection: EDM/rigidity is
  intrinsically invariant to rigid transforms, so \citep{iiyama2024,iiyama2025edm}
  \emph{exploit} the spatial datum gauge as a tool---that is why they need no precise
  ephemeris---whereas we \emph{characterize} that same invariance as the fault-hiding subspace
  and extend it to the timescale.
  \item \textbf{Versus secure state estimation \citep{fawzi2014,shoukry2016}.} The
  sparse-observability / $2f$-bound is theirs. Our Theorem~\ref{thm:byzantine} is its
  \emph{geometric instantiation} on the coupled-datum lunar Jacobian $\Gmat$, with the
  peer-column-block structure of a physical constellation; it is not a new bound.
  \item \textbf{Versus LiAISON \citep{hillborn2007}.} LiAISON recovers absolute state from
  crosslinks by exploiting \emph{orbital dynamics} over an arc---a dynamically-unique orbit
  breaks the frame symmetry. Our analysis is the \emph{kinematic / snapshot} regime: the
  instantaneous geometry Jacobian, whose gauge is eight. The dynamical regime is
  complementary and can partially collapse that gauge; we state this scope explicitly
  (Section~\ref{sec:model}) rather than overclaim a dynamics-free absolute frame.
\end{itemize}
What is genuinely ours is therefore narrow and, we hope, useful: (i)~the coupled
frame$\oplus$timescale gauge instantiated in fault observability---the undetectable-fault
subspace is the measurement-space \emph{face} of the same rank deficiency that produces the
datum$\oplus$timescale gauge in parameter space; (ii)~the
geometric secure-\emph{identification} bound for colluding peers on that geometry; (iii)~the
unification tying P1/P3/P4 to fault-hiding; and (iv)~the provider-fault taxonomy and its
common/differential detectability split.

\paragraph{Contributions.}
\begin{enumerate}
  \item \textbf{The one-rank-deficiency unification} (\Proven{}, Section~\ref{sec:obs}). The
  same rank deficiency of $\Gmat$ produces the parameter-space gauge
  $\Nsp(\Gmat)\subset\R^{40}$ and the measurement-space blind fault subspace
  $\rng(\Gmat)\subset\R^{84}$; an undetectable fault is one absorbed into a state error
  defined only modulo $\Nsp(\Gmat)$. This is a statement about \emph{two subspaces linked by
  one map}, not an identity of subspaces.
  \item \textbf{Undetectable $\iff\rng(\Gmat)$, surviving the gauge} (\Cited{} core,
  \Proven{} survival, Theorem~\ref{thm:undetectable}). The classical parity characterization
  holds for the rank-deficient, gauge-carrying $\Gmat$ through the pseudo-inverse projector,
  which is gauge-independent.
  \item \textbf{A geometric secure-identification bound} (\Cited{} bound, \Proven{}
  instantiation, Theorem~\ref{thm:byzantine}). Detect $\le f$ colluding peers iff the
  parity-projected peer blocks are full column rank (block-spark $>f$); identify iff
  block-spark $>2f$. The condition is a measurement-redundancy coding bound, distinct from
  the $3f+1$ consensus threshold.
  \item \textbf{An autonomous-holdover floor and a provider taxonomy} (\Proven{} structure,
  \Modelled{} magnitude, Section~\ref{sec:holdover}). The common clock rate is an irreducible
  temporal-gauge free parameter; a provider-wide bias splits into an undetectable common part
  (needs an external tie) and a detectable differential part (three or more providers
  identify the culprit).
  \item \textbf{The exactly-one validated pipeline} (\Validated{}, Section~\ref{sec:anchor})
  and \textbf{four design laws} (Section~\ref{sec:designlaws}), with a stated firewall and no
  flight, certification, or technology-readiness claim.
\end{enumerate}

\paragraph{Honesty scope.}
We separate four kinds of statement and label them at every occurrence. The parity/MDB/PL
machinery, the sparse-observability bound, the Byzantine model, and the rigidity and LiAISON
prior art are \emph{cited}. The gauge-survival of the parity characterization, the
one-rank-deficiency unification, the geometric block-spark instantiation, the holdover
temporal-gauge floor, the provider common/differential split, and the protection-gap
statement are \Proven{} algebra. Exactly one row is externally \Validated{}: the
linear-algebra pipeline reproduces an independent numpy/scipy oracle on real
DE440-orientation rows. Everything else---the representative geometry, every magnitude and
count read as an \emph{operational} tolerance, PL, provider-count, or holdover drift---is
\Modelled{}. There is no flight data, no certified alert limit, no endorsement by any agency,
and no technology-readiness claim above analysis; every number is regenerated from the open
engine by one command (Section~\ref{sec:availability}).

\section{The self-referential network model}
\label{sec:model}

We work with a linearized model in which a sparse lunar constellation---orbiters and
surface nodes from more than one provider---estimates its own state from the ranges and
clock ties its nodes form with one another, with \emph{no} node position or time known
externally. This is the free (self-referential) network, the setting P4 identified as the
carrier of the coupled datum$\oplus$timescale gauge \citep{baweja2026gauge}; here it is the
setting in which an autonomous integrity monitor must work. This section fixes the state,
the observation model, the two subspaces that the rest of the paper keeps rigidly apart, the
fault taxonomy, and the kinematic scope. All quantities named \code{like\_this} are the
corresponding items in the open engine module \code{lunar\_faultobs}
(Section~\ref{sec:availability}).

\subsection{Nodes, state, and the per-node coordinates}
\label{sec:state}

The network has $M=8$ nodes: five near-side surface nodes (surface transponder beacons) and
three relay orbiters, drawn from at least two providers. Each node $j$ carries a local
five-parameter deviation coordinate
\begin{equation}
  \xvec_j \;=\; \bigl[\,\underbrace{\delta\bm p_j}_{\in\,\R^3},\;
                       \underbrace{\dtau_j}_{\text{offset}},\;
                       \underbrace{\dalpha_j}_{\text{rate}}\,\bigr]\;\in\;\R^{5},
  \qquad
  \xvec \;=\; [\,\xvec_1;\dots;\xvec_M\,]\;\in\;\R^{5M}=\R^{40},
  \label{eq:state}
\end{equation}
a position deviation $\delta\bm p_j$ (metres) and a clock offset $\dtau_j$ and rate
$\dalpha_j$. With $M=8$ the state dimension is $5M=40$; in the engine the layout constants
are \code{PARAMS\_PER\_NODE}${}=5$, \code{state\_dim}${}=40$.

\subsection{The differential observation model}
\label{sec:obsmodel}

Observations are inter-node \emph{differential} one-way pseudoranges and clock ties. A
differential row is built so that a whole-network Helmert transform (translation, rotation)
\emph{and} a common clock offset and rate cancel identically---both endpoints of an
inter-node measurement move together---so the row is invariant under the network's own datum
freedom. The engine builders \code{differential\_range\_row} and
\code{differential\_clock\_tie\_row} form these rows (the free-network analogue of P4's
anchored \code{oneway\_range\_row}); the spatial partials are computed through the
\emph{real} DE440 lunar physical-axis orientation path, so the line-of-sight geometry is the
external ingredient, while the node states are \emph{estimated}, not pinned.

Each of the $\binom{M}{2}=\binom{8}{2}=28$ node pairs is observed at three clock-integration
windows $T\in\{21600,43200,86400\}$~s (six hours, twelve hours, one day), giving
\begin{equation}
  n \;=\; \binom{8}{2}\times 3 \;=\; 28\times 3 \;=\; 84
  \label{eq:nrows}
\end{equation}
rows. Stacking the rows gives the linearized model
\begin{equation}
  \bm z \;=\; \Gmat\,\xvec \;+\; \bvec \;+\; \bm\varepsilon,
  \qquad
  \bm\varepsilon\sim\mathcal N(\bm 0,\Sig),
  \qquad
  \Wmat \;:=\; \Sig^{-1}\;=\;\operatorname{diag}(1/\sigma_i^2),
  \label{eq:model}
\end{equation}
with $\Gmat\in\R^{84\times 40}$ the network Jacobian, $\bvec\in\R^{84}$ a fault (bias)
vector, and $\Wmat$ the weight matrix. The measurement standard deviations cycle through
$\sigma\in\{0.5,1,2\}$~m, so $\Wmat\neq\bm{I}$ is a genuinely non-uniform weighting---which,
as Section~\ref{sec:obs} shows, is exactly the regime in which the correct minimum
detectable bias form matters. The epoch is 2024-01-01 TT (\code{t\_tt\_jc}${}=0.239986$
Julian centuries), fixing the DE440 orientation.

\subsection{The gauge: what the self-referential network cannot observe}
\label{sec:gauge}

Because no node is externally pinned, the network Jacobian is rank-deficient. Its null space
is the datum$\oplus$timescale \emph{gauge} of P4 \citep{baweja2026gauge}, realized
differentially:
\begin{equation}
  \Nsp(\Gmat)
  \;=\;\operatorname{span}\{\,
    \underbrace{3\text{ translations}}_{\text{rigid}},\;
    \underbrace{3\text{ rotations}}_{\text{rigid}},\;
    \underbrace{\text{common }\dtau}_{\text{offset}},\;
    \underbrace{\text{common }\dalpha}_{\text{rate}}\,\},
  \qquad \dim\Nsp(\Gmat)=8 .
  \label{eq:gauge}
\end{equation}
A common translation and infinitesimal rotation of the whole point set, a common clock
offset, and a common clock rate leave every internal differential observation unchanged, so
the eight corresponding directions are annihilated by $\Gmat$. Scale is \emph{observable},
not a gauge direction, because under the SI second and the $c$-defined metre a differential
range carries absolute length (the P4 result, imported unchanged). Since $\Wmat\succ0$,
$\Nsp(\Gmat^\top\Wmat\Gmat)=\Nsp(\Gmat)$, so the weighted normal matrix has the same
eight-dimensional gauge. With $\dim\Nsp(\Gmat)=8$ and state dimension $40$, the rank is
$\rank\Gmat=40-8=32$.

\subsection{Two subspaces, kept distinct}
\label{sec:twosubspaces}

The whole paper turns on holding two subspaces apart; we name them once, here.
\begin{itemize}
  \item The \textbf{parameter-space gauge}
  $\Nsp(\Gmat)\subset\R^{40}$ of \eqref{eq:gauge}: the unobservable \emph{state} moves---the
  absolute frame (P1) and the common clock offset and rate (P4). It lives in the state space
  $\R^{5M}$.
  \item The \textbf{measurement-space blind fault subspace}
  $\rng(\Gmat)\subset\R^{84}$: the \emph{biases} that leave the parity residual zero and are
  therefore undetectable (Section~\ref{sec:obs}). The common-mode error class of P3
  \citep{baweja2026commonmode} is a subset of this. It lives in the measurement space
  $\R^{n}$.
\end{itemize}
These are different subspaces of different spaces, with different dimensions ($8$ versus
$32$). What links them is that they are the two faces of the \emph{same} rank deficiency of
$\Gmat$, and Section~\ref{sec:obs} makes the link precise: an undetectable fault in
$\rng(\Gmat)$ is absorbed into a state error defined only modulo $\Nsp(\Gmat)$. We never
write that the two subspaces are ``the same.''

\subsection{Fault taxonomy}
\label{sec:taxonomy}

A fault is a bias vector $\bvec\in\R^{84}$; its \emph{support} distinguishes three classes,
of increasing adversarial power.
\begin{description}
  \item[Measurement fault.] A bias on a single measurement row: $\bvec=\gamma\,\bm e_i$ for
  one row $i$. This is the classical RAIM single-fault model.
  \item[Peer (Byzantine) fault.] A node $j$ corrupts, arbitrarily and possibly in collusion
  with other nodes, \emph{all} the measurements incident to it. Its signature is the
  column-block $\Bmat_j\in\R^{84\times m_j}$ selecting node $j$'s $m_j$ incident rows; the
  injected bias is $\bvec=\Bmat_j\bm v$ for arbitrary $\bm v\in\R^{m_j}$. In our network a
  node participates in $7$ pairs $\times\,3$ windows $=21$ rows, so a broad-coverage peer
  owning all of them has $m_j=21$. This is the Byzantine model of \citep{lamport1982},
  used for the fault \emph{model} only.
  \item[Provider fault.] A provider-wide ephemeris or timescale error is a \emph{coordinated}
  bias across all of that provider's nodes---a structured, multi-node fault treated in
  Section~\ref{sec:holdover}.
\end{description}

\subsection{Kinematic (snapshot) scope}
\label{sec:scope}

$\Gmat$ is the \emph{instantaneous geometry} Jacobian at the epoch: the gauge \eqref{eq:gauge}
of eight is the kinematic / snapshot datum defect. This is the scope of the paper, and it is
a real limitation. LiAISON \citep{hillborn2007} shows that \emph{orbital dynamics} over an
arc can recover absolute state from crosslinks by exploiting a dynamically-unique orbit; such
dynamical information can partially collapse $\Nsp(\Gmat)$. Our results are the geometry /
snapshot regime; the dynamical regime is complementary and is not claimed here. The three
integration windows in \eqref{eq:nrows} carry \emph{relative} clock information (per-node
rates become observable across windows), but not the \emph{common} rate, which remains a
gauge direction (Section~\ref{sec:holdover})---consistent with the kinematic scope.

\section{Fault observability and the one rank deficiency}
\label{sec:obs}

This section states the paper's structural core: the parity projector that survives the
rank-deficient, gauge-carrying $\Gmat$; the characterization of undetectable faults as
exactly $\rng(\Gmat)$; the one-rank-deficiency unification that ties the parameter-space
gauge to the measurement-space blind subspace; and the protection-gap consequence. The
classical parity result is \Cited{}; the gauge-survival, the unification, and the
protection gap are \Proven{}; the engine numbers that instantiate them are outputs of the
\Validated{} pipeline of Section~\ref{sec:anchor}.

\subsection{The weighted parity projector}
\label{sec:parity}

The weighted least-squares estimate of $\xvec$ from \eqref{eq:model} is
$\hat{\xvec}=(\Gmat^\top\Wmat\Gmat)^{+}\Gmat^\top\Wmat\bm z$, using the Moore--Penrose
pseudo-inverse because $\Gmat^\top\Wmat\Gmat$ is rank-deficient (its kernel is the gauge
\eqref{eq:gauge}). The residual is $\bm r=\Pperp\bm z$ with the \emph{weighted parity
projector}
\begin{equation}
  \Pperp \;=\; \bm{I} \;-\; \Gmat\,(\Gmat^\top\Wmat\Gmat)^{+}\,\Gmat^\top\Wmat
  \;\in\;\R^{84\times84}.
  \label{eq:pperp}
\end{equation}

\begin{proposition}[The parity projector is gauge-independent]
\label{prop:projector}
Let $\bm M:=\Wmat^{1/2}\Gmat$ and let $\Pi_{\bm M}$ be the orthogonal projector onto
$\rng(\bm M)$. Then $\Pperp=\Wmat^{-1/2}(\bm{I}-\Pi_{\bm M})\Wmat^{1/2}$ is idempotent and
$\Wmat$-self-adjoint (i.e.\ $\Wmat\Pperp=\Pperp{}^{\!\top}\Wmat$, \emph{not} ordinary
symmetric when $\Wmat\neq\bm{I}$), with $\Pperp\Gmat=\bm0$, $\ker\Pperp=\rng(\Gmat)$, and
$\rng(\Pperp)$ equal to the $\Wmat$-orthogonal complement of $\rng(\Gmat)$. Every one of
these objects depends only on $\rng(\Gmat)$ and $\Wmat$, hence is invariant under the
gauge: adding any $\bm n\in\Nsp(\Gmat)$ to the state leaves $\Pperp$ and $\bm r$ unchanged.
\end{proposition}

\begin{proof}
$\Gmat(\Gmat^\top\Wmat\Gmat)^{+}\Gmat^\top\Wmat$ is the $\Wmat$-orthogonal (oblique in the
standard inner product) projector onto $\rng(\Gmat)$: writing $\bm M=\Wmat^{1/2}\Gmat$, one
has $\Wmat^{1/2}\Gmat(\Gmat^\top\Wmat\Gmat)^{+}\Gmat^\top\Wmat^{1/2}=\bm M(\bm M^\top\bm
M)^{+}\bm M^\top=\Pi_{\bm M}$, the ordinary orthogonal projector onto $\rng(\bm M)$, whose
defining properties (idempotence, self-adjointness, range) are basis-free and survive the
rank deficiency because $(\bm M^\top\bm M)^{+}$ is the pseudo-inverse. Conjugating by
$\Wmat^{-1/2}$ gives \eqref{eq:pperp}, so $\Pperp=\Wmat^{-1/2}(\bm{I}-\Pi_{\bm M})\Wmat^{1/2}$
is idempotent, and $\Wmat\Pperp=\Wmat^{1/2}(\bm{I}-\Pi_{\bm M})\Wmat^{1/2}$ is symmetric,
i.e.\ $\Pperp$ is $\Wmat$-self-adjoint. Finally $\Pperp\Gmat=\Wmat^{-1/2}(\bm{I}-\Pi_{\bm
M})\bm M=\bm0$, so $\rng(\Gmat)\subseteq\ker\Pperp$; equality follows by dimension count,
$\dim\ker\Pperp=\rank\Gmat=\dim\rng(\Gmat)$. None of $\rng(\Gmat)$, $\Wmat$, or $\Pi_{\bm
M}$ references a basis of $\Nsp(\Gmat)$, so all are gauge-independent.
\end{proof}

Because $\mathbb E[\bm r]=\Pperp\,\mathbb E[\bm z]=\Pperp\bvec$ (the signal $\Gmat\xvec$ is
annihilated by $\Pperp$), the gauge does \emph{not} leak into detectability: the parity
residual sees a fault through $\Pperp$ alone.

\subsection{Undetectable faults are exactly \texorpdfstring{$\rng(\Gmat)$}{range(G)}}
\label{sec:undetectable}

\begin{theorem}[Undetectable-fault characterization, gauge-carrying $\Gmat$]
\label{thm:undetectable}
A fault $\bvec$ produces zero expected parity residual, $\mathbb E[\bm r]=\Pperp\bvec=\bm0$,
if and only if $\bvec\in\rng(\Gmat)$. The characterization holds for the rank-deficient,
gauge-carrying free-network $\Gmat$ through the pseudo-inverse projector \eqref{eq:pperp}.
Consequently, in a network with no redundancy ($n=\rank\Gmat$) one has $\rng(\Gmat)=\R^{n}$
and \emph{nothing} is detectable.
\end{theorem}

\begin{proof}
By Proposition~\ref{prop:projector}, $\ker\Pperp=\rng(\Gmat)$, so $\Pperp\bvec=\bm0\iff
\bvec\in\rng(\Gmat)$. The classical parity statement ``$\bvec$ undetectable
$\iff\bvec\in\rng(\Gmat)$'' is \Cited{} \citep{sturza1988,brown1992,parkinson1988}; the only
new content is that it survives a rank-deficient $\Gmat$, which Proposition~\ref{prop:projector}
supplies by using $(\Gmat^\top\Wmat\Gmat)^{+}$ in place of an inverse and by noting $\Pperp$
is gauge-independent. If $n=\rank\Gmat$ then $\rng(\Gmat)=\R^{n}$, so every $\bvec$ lies in
$\rng(\Gmat)$ and $\Pperp=\bm0$.
\end{proof}

The pipeline evaluates two witnesses on the real-DE440 network (both counts \Validated{},
Section~\ref{sec:anchor}). An \emph{in-range} fault $\bvec=\Gmat\,\bm e_{\dtau(\text{node
}1)}$---a shift of node~1's clock offset---has parity norm $\lVert\Pperp\bvec\rVert=9.581
\times10^{-12}$ and is \textbf{undetectable}. A \emph{generic} fault $\bvec=\bm e_0$ (a
single measurement bias) has $\lVert\Pperp\bvec\rVert=1.021$ and is \textbf{detectable}.

\subsection{Minimum detectable bias}
\label{sec:mdb}

For a unit fault direction $\bm c$ (a column of a fault signature), the Baarda minimum
detectable bias is the smallest fault magnitude a non-central parity test rejects at
prescribed false-alarm and detection powers \citep{baarda1968,teunissen2006}:
\begin{equation}
  \MDB(\bm c)\;=\;\sqrt{\frac{\lambda_0}{\bm c^\top\Wmat\Pperp\bm c}},
  \label{eq:mdb}
\end{equation}
where $\lambda_0$ is the non-centrality set by the false-alarm probability $P_{\mathrm{fa}}$
and the missed-detection probability $P_{\mathrm{md}}$. The non-centrality is the
\emph{single}-projector, $\Wmat$-weighted parity energy $\bm c^\top\Wmat\Pperp\bm c$; because
$\Pperp$ is $\Wmat$-self-adjoint and idempotent (Proposition~\ref{prop:projector}), it equals
$(\Pperp\bm c)^\top\Wmat(\Pperp\bm c)$ for \emph{all} $\Wmat$---so this correct form is the
weighted energy of the projected column. It is \emph{not} the double-projector form
$\bm c^\top\Pperp\Wmat\Pperp\bm c$, which differs whenever $\Wmat\neq\bm{I}$ and is the wrong
non-centrality; Appendix~\ref{app:mdb} states the identity and the counterexample precisely.
With $\lambda_0=17.075$ from $P_{\mathrm{fa}}=10^{-3}$, $P_{\mathrm{md}}=0.20$, the pipeline
reports (both \Validated{}) $\MDB(\bm e_0)=4.6592$ for the generic single-measurement fault,
and $\MDB=\infty$ for the in-range fault $\Gmat\,\bm e_{\dtau(\text{node }1)}$, whose
non-centrality $\bm c^\top\Wmat\Pperp\bm c=0$---the protection gap of
Proposition~\ref{prop:protection} read as a bias magnitude.

\subsection{The one rank deficiency, two faces}
\label{sec:unification}

\begin{proposition}[One rank deficiency, two faces]
\label{prop:unification}
The single rank deficiency of the self-referential Jacobian, $\rank\Gmat=32<40$, produces
\emph{both}
\begin{enumerate}
  \item[(a)] the parameter-space gauge $\Nsp(\Gmat)\subset\R^{40}$, of dimension
  $40-\rank\Gmat=8$ (the datum$\oplus$timescale gauge, P1$\,\oplus\,$P4); and
  \item[(b)] the measurement-space blind fault subspace $\rng(\Gmat)\subset\R^{84}$, of
  dimension $\rank\Gmat=32$ (the undetectable faults of Theorem~\ref{thm:undetectable}, with
  the P3 common-mode class a subset).
\end{enumerate}
The two are linked, not equal: for any undetectable fault $\bvec\in\rng(\Gmat)$ there is a
state shift $\Delta\xvec$ with $\Gmat\,\Delta\xvec=\bvec$, and $\Delta\xvec$ is determined
only \emph{modulo} $\Nsp(\Gmat)$. Thus an undetectable fault is precisely one absorbed into
a state error whose ambiguity is exactly the gauge.
\end{proposition}

\begin{proof}
The rank--nullity theorem applied to $\Gmat\colon\R^{40}\to\R^{84}$ gives
$\dim\Nsp(\Gmat)=40-\rank\Gmat=8$ and $\dim\rng(\Gmat)=\rank\Gmat=32$; both are functions of
the one integer $\rank\Gmat$, which is (a) and (b). For the link, $\bvec\in\rng(\Gmat)$ means
$\bvec=\Gmat\,\Delta\xvec$ has a solution by definition of the range; the solution set is the
coset $\Delta\xvec+\Nsp(\Gmat)$, since $\Gmat(\Delta\xvec+\bm n)=\Gmat\Delta\xvec$ for
$\bm n\in\Nsp(\Gmat)$ and only for such $\bm n$. Hence the induced state error is defined
modulo $\Nsp(\Gmat)$.
\end{proof}

This is the honest form of the program's unification: it is a statement about two subspaces
tied by one map, \emph{not} the (false) claim that the gauge and the blind subspace are the
same set. The pipeline reports the two dimensions directly---$\rank\Gmat=32$ and the parity
dimension $\operatorname{trace}(\Pperp)=84-32=52$ (both \Validated{}), against the eight
gauge directions $\Nsp(\Gmat)$ that Theorem P4 establishes (\Proven{}). The eight analytic
gauge generators are annihilated by $\Gmat$ to a worst relative residual of $2.702\times
10^{-17}$, confirming $\Nsp(\Gmat)\supseteq\operatorname{span}\{\text{3 transl},\text{3
rot},\dtau,\dalpha\}$ exactly. The count is spectrally unambiguous: the largest
gauge-direction eigenvalue of the (preconditioned) normal matrix is
$\lambda_{\text{null}}^{\max}=4.367\times10^{-15}$, the smallest observable eigenvalue is
$\lambda_{\text{obs}}^{\min}=5.455\times10^{-4}$, and the largest is
$\lambda^{\max}=8.597\times10^{1}$, a null-to-observable spectral gap of $1.249\times
10^{11}$.

\begin{remark}[Common-mode is a subset of the blind subspace]
\label{rem:commonmode}
The common-mode reference-frame error class of P3 \citep{baweja2026commonmode}---a shift
$\delta\bm y=\Gmat[\,-\Delta\bm s;\,\bm 0\,]$ that moves the frame realization---is of the
form $\Gmat\,(\cdot)$ and so lies in $\rng(\Gmat)$ by construction; it is a special case of
Theorem~\ref{thm:undetectable}, undetectable to any snapshot residual test. Paired with the
parameter-space gauge $\Nsp(\Gmat)$ (P1/P4), this is the two-faces statement of
Proposition~\ref{prop:unification}: the frame that a monitor cannot pin (parameter space)
and the frame-shaped fault it cannot see (measurement space) are the two images of the one
rank deficiency.
\end{remark}

\subsection{The protection gap}
\label{sec:protectiongap}

Theorem~\ref{thm:undetectable} has a protection-level consequence. For an undetectable fault
the residual is zero but the state estimate still moves; the ratio of the two is the Brown
slope, and for the $\rng(\Gmat)$ class it is infinite.

\begin{proposition}[The protection gap: infinite slope on $\rng(\Gmat)$]
\label{prop:protection}
Let $\Pi_{\mathrm{obs}}$ project the state error onto the observable (user-PVT) directions,
and define the weighted Brown slope of a fault as
\begin{equation}
  \operatorname{slope}(\bvec)
  \;=\;\frac{\lVert\Pi_{\mathrm{obs}}\,\Delta\hat{\xvec}\rVert}{\lVert\Pperp\bvec\rVert_{\Wmat}},
  \qquad
  \Delta\hat{\xvec}=(\Gmat^\top\Wmat\Gmat)^{+}\Gmat^\top\Wmat\,\bvec .
  \label{eq:slope}
\end{equation}
For $\bvec\in\rng(\Gmat)$ with $\Pi_{\mathrm{obs}}\Delta\hat{\xvec}\neq\bm0$ the denominator
vanishes ($\lVert\Pperp\bvec\rVert_{\Wmat}=0$ by Theorem~\ref{thm:undetectable}) while the
numerator does not, so $\operatorname{slope}(\bvec)=\infty$: the fault corrupts an observable
state combination while producing no parity residual. No finite alert limit protects this
fault class autonomously.
\end{proposition}

\begin{proof}
By Theorem~\ref{thm:undetectable}, $\bvec\in\rng(\Gmat)\Rightarrow\Pperp\bvec=\bm0\Rightarrow
\lVert\Pperp\bvec\rVert_{\Wmat}=0$. The estimator response
$\Delta\hat{\xvec}=(\Gmat^\top\Wmat\Gmat)^{+}\Gmat^\top\Wmat\bvec$ need not lie in
$\Nsp(\Gmat)$, so its observable projection $\Pi_{\mathrm{obs}}\Delta\hat{\xvec}$ can be
nonzero; the in-range witness of Section~\ref{sec:undetectable} (a node-1 offset shift, an
\emph{observable} relative-clock combination) is such a case. A ratio with zero denominator
and nonzero numerator is $+\infty$.
\end{proof}

The finite counterpart---the worst \emph{detectable} single-measurement slope and the
alert-limit crossing---is a \Modelled{} magnitude quantified in Section~\ref{sec:designlaws}.

\section{Byzantine integrity: a geometric secure-identification bound}
\label{sec:byzantine}

An autonomous constellation cannot assume its peers are honest. A compromised or
malfunctioning node may inject \emph{arbitrary}, possibly colluding, biases on all the
measurements it participates in. This section asks how many such Byzantine peers the network
can detect and identify from its internal geometry alone. The answer is a
\emph{secure-identification} bound: a block-structured instance of the sparse-observability /
error-correcting-code bound of secure state estimation \citep{fawzi2014,shoukry2016,%
candes2005,donoho2003}. The bound is \Cited{}; its instantiation on the coupled-datum lunar
$\Gmat$ is \Proven{}; the tolerance \emph{counts} are \Validated{} pipeline outputs, and any
reading of them as an operational flight guarantee is \Modelled{}.

\subsection{Peer signatures and the effective (parity-projected) block}
\label{sec:peersig}

A peer fault at node $j$ has signature $\Bmat_j\in\R^{84\times m_j}$, the column-block
selecting node $j$'s $m_j$ incident measurement rows (Section~\ref{sec:taxonomy}); the
injected bias is $\bvec=\Bmat_j\bm v$ for an arbitrary $\bm v\in\R^{m_j}$. A coalition
$T\subseteq\{1,\dots,M\}$ of Byzantine peers injects $\bvec=\Bmat_T\bm v_T$, where
$\Bmat_T=[\Bmat_j]_{j\in T}$ stacks the coalition's column-blocks. What the parity test sees
is the \emph{effective signature}
\begin{equation}
  \Gbar_j \;:=\; \Pperp\Bmat_j,
  \qquad
  \Gbar_T \;=\; \Pperp\Bmat_T ,
  \label{eq:effsig}
\end{equation}
the peer block after the parity projector \eqref{eq:pperp}. A coalition is invisible exactly
when its injected bias can be hidden in $\rng(\Gmat)$, i.e.\ when $\Gbar_T$ is column-rank
deficient.

\begin{definition}[Block-spark]
\label{def:blockspark}
The \emph{block-spark} of the family $\{\Gbar_j\}_{j=1}^{M}$ is the smallest number of
peer-blocks whose union is column-rank deficient:
\begin{equation}
  \spark_{\mathrm b}\bigl(\{\Gbar_j\}\bigr)
  \;=\;\min\Bigl\{\,|T|\;:\;\Gbar_T=\Pperp\Bmat_T\ \text{is column-rank deficient}\,\Bigr\}.
  \label{eq:blockspark}
\end{equation}
It is the block analogue of the spark of a dictionary \citep{donoho2003} (the smallest number
of columns that are linearly dependent), lifted from single columns to peer column-blocks.
\end{definition}

\subsection{The detect / identify bound}
\label{sec:bound}

\begin{theorem}[Geometric secure detect / identify bound]
\label{thm:byzantine}
Let $\{\Gbar_j\}$ be the effective peer signatures \eqref{eq:effsig} on the free-network
$\Gmat$. Then:
\begin{enumerate}
  \item[(i)] \textbf{Detection.} All coalitions of size $\le f$ are detectable if and only if
  $\Gbar_T$ is full column rank for every $|T|\le f$, equivalently
  $\spark_{\mathrm b}(\{\Gbar_j\})>f$, equivalently $\operatorname{span}(\Bmat_T)\cap
  \rng(\Gmat)=\{\bm0\}$ for every $|T|\le f$.
  \item[(ii)] \textbf{Identification.} All coalitions of size $\le f$ are uniquely
  identifiable if and only if $\Gbar_U$ is full column rank for every $|U|\le 2f$,
  equivalently $\spark_{\mathrm b}(\{\Gbar_j\})>2f$ (a coalition minimum distance $\ge2f+1$).
\end{enumerate}
\end{theorem}

\begin{proof}
\emph{(i).} A size-$\le f$ coalition $T$ is undetectable iff there is a nonzero injection
$\Bmat_T\bm v$ with $\Pperp\Bmat_T\bm v=\bm0$, i.e.\ iff $\Gbar_T=\Pperp\Bmat_T$ has a
nontrivial kernel, i.e.\ iff $\Gbar_T$ is column-rank deficient. So \emph{all}
$|T|\le f$ are detectable iff every such $\Gbar_T$ is full column rank, which by
Definition~\ref{def:blockspark} is $\spark_{\mathrm b}>f$. By
Proposition~\ref{prop:projector}, $\Pperp\Bmat_T\bm v=\bm0\iff\Bmat_T\bm v\in\ker\Pperp=
\rng(\Gmat)$, giving the intersection form. \emph{(ii).} Two coalitions $T_1,T_2$ with
$|T_i|\le f$ produce the same residual iff $\Pperp(\Bmat_{T_1}\bm v_1-\Bmat_{T_2}\bm v_2)=
\bm0$, i.e.\ iff some injection on $T_1\cup T_2$ ($|T_1\cup T_2|\le2f$) is undetectable.
Unique identification of every $|T|\le f$ therefore requires every $|U|\le2f$ union to be
detectable, i.e.\ $\Gbar_U$ full column rank for all $|U|\le2f$, i.e.\ $\spark_{\mathrm
b}>2f$. This is the block form of the classical $2f$-sparse-observability / minimum-distance
$\ge2f+1$ condition of \citep{fawzi2014,shoukry2016,candes2005}.
\end{proof}

\paragraph{The detect condition is full column rank, not ``$\Pperp\Bmat_T\neq\bm0$.''}
It is tempting to declare a coalition detectable as soon as its projected block is nonzero.
This is \emph{wrong}. A block $\Gbar_T=\Pperp\Bmat_T$ can be nonzero yet column-rank
deficient, and the adversary then hides a fault in the null combination---an injection
$\Bmat_T\bm v\in\rng(\Gmat)$ that produces \emph{zero} residual even though $\Gbar_T\neq
\bm0$ on other injections. The correct detect condition is that every $\le f$-block union be
\emph{full column rank} (span disjoint from $\rng(\Gmat)$); the identify condition raises
this to $2f$.

\subsection{The Lamport distinction: this is not a \texorpdfstring{$3f+1$}{3f+1} consensus bound}
\label{sec:lamport}

We adopt the Byzantine fault \emph{model}---arbitrary, colluding, worst-case node
behaviour---from Lamport--Shostak--Pease \citep{lamport1982}, and \emph{only} the model. The
bound of Theorem~\ref{thm:byzantine} is a \emph{measurement-redundancy / coding} bound: it
counts how many arbitrarily corrupted \emph{measurement blocks} a geometric estimator can
detect ($\spark_{\mathrm b}>f$, i.e.\ $f+1$ block-redundancy, for detection) and correct
($\spark_{\mathrm b}>2f$ for identification), exactly as in secure state estimation and
error-correcting codes
\citep{fawzi2014,shoukry2016,candes2005,donoho2003}. It is \emph{not}, and must not be read
as, the $3f+1$ threshold for Byzantine agreement in a message-passing consensus network
\citep{lamport1982}. Those are different problem classes: consensus concerns how many
faulty \emph{processes} exchanging \emph{messages} still permit agreement; ours concerns how
many faulty \emph{measurement blocks} in a \emph{static estimation} geometry still permit
fault correction. The two thresholds ($3f+1$ versus $\spark_{\mathrm b}>2f$) are not
comparable, and we draw no analogy between them.

\subsection{Redundancy raises tolerance; broad coverage destroys it}
\label{sec:counts}

The bound of Theorem~\ref{thm:byzantine} is governed by the network's block-spark, and the
pipeline computes it by exact enumeration for the representative geometry (the counts are
\Validated{}; their reading as an operational guarantee is \Modelled{}).

\paragraph{Redundant network.}
A network of six independent single-measurement peers---each owning one linearly independent
row---has $\spark_{\mathrm b}=7$, giving $f_{\mathrm{detect}}=6$ and $f_{\mathrm{identify}}=3$
(the largest $f$ with $\spark_{\mathrm b}>2f$). Redundancy buys tolerance: adding independent
cross-links or providers raises the block-spark and hence $f$.

\paragraph{Coverage fragility and the deficiency-of-five self-blindness.}
The opposite regime is instructive and is the paper's cautionary result. A \emph{broad}-%
coverage peer that owns \emph{all} of node~2's incident rows (all $21$ of them) has
$\spark_{\mathrm b}=1$ and $f_{\mathrm{detect}}=0$: a \emph{single} such peer is already
\textbf{undetectable}. A \emph{narrow} peer owning one measurement has $\spark_{\mathrm b}=2$
and $f_{\mathrm{detect}}=1$, and is detectable. The mechanism is precise: a peer owning all
of a node's measurements can inject a \emph{uniform self-bias} equal to a shift of that
node's own five local degrees of freedom---its three position components, its clock offset,
and its clock rate. That injection is of the form $\Gmat(\text{shift of node's local
block})\in\rng(\Gmat)$, hence undetectable by Theorem~\ref{thm:undetectable}: \emph{a node is
blind to a fault in its own five local DOF}. Broad measurement coverage by a single peer is
therefore the fragility, not the strength. This is why the detect condition is full column
rank of the union and not mere nonzeroness: the broad-coverage block is nonzero after
projection on most injections, yet column-rank deficient on the self-bias direction.

\begin{table}[t]
\centering
\caption{Byzantine block-spark on the representative network (real DE440 orientation).
Counts are \Validated{} (matched exactly by the independent oracle, Section~\ref{sec:anchor});
their reading as an operational flight tolerance is \Modelled{}. $f_{\mathrm{detect}}=
\spark_{\mathrm b}-1$; $f_{\mathrm{identify}}=\lfloor(\spark_{\mathrm b}-1)/2\rfloor$.}
\label{tab:byzantine}
\small
\setlength{\tabcolsep}{7pt}
\begin{tabular}{@{}lccc@{}}
\toprule
\textbf{Peer structure} & $\spark_{\mathrm b}$ & $f_{\mathrm{detect}}$ & $f_{\mathrm{identify}}$ \\
\midrule
Six independent single-measurement peers & $7$ & $6$ & $3$ \\
Narrow peer (one incident measurement)   & $2$ & $1$ & $0$ \\
Broad peer (all $21$ rows of node 2)     & $1$ & $0$ & $0$ \\
\bottomrule
\end{tabular}
\end{table}

\subsection{Provider faults: the common / differential split}
\label{sec:providersplit}

The most structured Byzantine fault is a \emph{provider}-wide bias---a coordinated error
across all nodes of one provider (Section~\ref{sec:taxonomy}). Such a fault splits into a
part all providers share and a part that distinguishes them.

\begin{proposition}[Provider common / differential detectability]
\label{prop:provider}
A provider-wide bias decomposes as $\bvec=\bvec_{\mathrm{common}}+\bvec_{\mathrm{diff}}$,
where the common part---shared identically by every provider---lies in $\rng(\Gmat)$ and is
therefore \emph{internally undetectable} (it is the common-mode class of
Remark~\ref{rem:commonmode}, requiring an external tie, and is the common component of the
inter-ephemeris floor of \citep{baweja2026interop}); and the differential part---provider $A$
disagreeing with provider $B$---has nonzero parity, $\Pperp\bvec_{\mathrm{diff}}\neq\bm0$,
and is \emph{internally detectable}. Consequently: with two providers the network
\emph{detects} disagreement but cannot \emph{identify} the culprit (the two-clock problem);
with three or more providers it identifies the culprit via block-spark $>2f$
(Theorem~\ref{thm:byzantine}).
\end{proposition}

\begin{proof}
The common part shifts every provider's realization by the same amount---a whole-network
datum move---so it is $\Gmat(\cdot)\in\rng(\Gmat)$ and $\Pperp\bvec_{\mathrm{common}}=\bm0$
by Theorem~\ref{thm:undetectable}. The differential part is supported on a proper subset of
provider blocks and is not a whole-network datum move, so it does not lie in $\rng(\Gmat)$
and has nonzero residual. Detection versus identification then follows from
Theorem~\ref{thm:byzantine}: two providers give two blocks whose disagreement is one
detectable direction but with no third block to break the tie ($\spark_{\mathrm b}=2$,
$f_{\mathrm{identify}}=0$); three or more give $\spark_{\mathrm b}>2$, enough to identify a
single faulty provider.
\end{proof}

The pipeline reports the number of internally detectable cross-provider rows as a function
of provider count (\Modelled{}): $1\to0$, $2\to48$, $4\to72$ (Table~\ref{tab:provider}). One
provider offers nothing to cross-check; two expose disagreement; four expose more. The count
grows sublinearly because the common part---the shared datum move---is never detectable no
matter how many providers share it.

\begin{table}[t]
\centering
\caption{Internally detectable cross-provider rows versus provider count (\Modelled{}). One
provider has no differential to expose; the shared common part stays in $\rng(\Gmat)$ and is
undetectable at any provider count (needs an external tie).}
\label{tab:provider}
\small
\setlength{\tabcolsep}{10pt}
\begin{tabular}{@{}ccl@{}}
\toprule
\textbf{Providers} & \textbf{Detectable cross-provider rows} & \textbf{Capability} \\
\midrule
$1$ & $0$  & no internal cross-check \\
$2$ & $48$ & detect disagreement, cannot identify culprit \\
$4$ & $72$ & detect and identify culprit (block-spark $>2f$) \\
\bottomrule
\end{tabular}
\end{table}

\section{Reproducibility: the validated linear-algebra pipeline}
\label{sec:anchor}

The paper's single externally \Validated{} claim is narrow and precise: that the
fault-observability \emph{linear-algebra pipeline}---the weighted parity projector
\eqref{eq:pperp}, the detectability norm $\lVert\Pperp\bvec\rVert$, the minimum detectable
bias \eqref{eq:mdb}, the rank and parity-dimension counts of
Proposition~\ref{prop:unification}, and the Byzantine block-spark of
Definition~\ref{def:blockspark}---reproduces an \emph{independent} numpy/scipy computation of
the same quantities, to relative and absolute error below $10^{-3}$ (integer counts exact),
on observation rows whose spatial partials carry the \emph{real} DE440 lunar physical-axis
orientation \citep{park2021de440}. This section states what is validated and, just as
carefully, what is not.

\subsection{The validated quantities}
\label{sec:validated}

The rows are built by \code{gen\_faultobs\_rows} through the real DE440 lunar physical-axis
orientation path, so every spatial column carries the physical libration at the fixture
epoch 2024-01-01 TT, not an approximation. On these rows the pipeline computes, and the
oracle reproduces, the quantities of Table~\ref{tab:anchor}.

\begin{table}[t]
\centering
\caption{The \Validated{} anchor: the fault-observability pipeline on real DE440-orientation
rows, reproduced by an independent numpy/scipy oracle to relative and absolute error below
$10^{-3}$ (integer counts exact). The projector-level maximum errors are $1.447\times
10^{-12}$ (absolute) and $6.6\times10^{-10}$ (relative), far inside the $10^{-3}$ tolerance;
the rank, parity-dimension, and block-spark counts match \emph{exactly}.}
\label{tab:anchor}
\footnotesize
\setlength{\tabcolsep}{6pt}
\begin{tabular}{@{}llll@{}}
\toprule
\textbf{Quantity} & \textbf{Value} & \textbf{Oracle check} & \textbf{Tag} \\
\midrule
Gauge dimension $\dim\Nsp(\Gmat)$          & $8$                    & exact & \Proven{} (P4) \\
Rank $\rank\Gmat$                          & $32$                   & exact & \Validated{} \\
Parity dimension $\operatorname{tr}\Pperp$ & $52$                   & exact & \Validated{} \\
$\lVert\Pperp\bvec\rVert$, in-range        & $9.581\times10^{-12}$  & $<10^{-3}$ & \Validated{} \\
$\lVert\Pperp\bvec\rVert$, generic $\bm e_0$& $1.021$               & $<10^{-3}$ & \Validated{} \\
$\MDB(\bm e_0)$ ($\lambda_0=17.075$)       & $4.6592$               & $<10^{-3}$ & \Validated{} \\
$\MDB$, in-range                           & $\infty$               & exact & \Validated{} \\
Block-spark, six single-meas peers         & $7$ ($f_{\mathrm d}=6$, $f_{\mathrm i}=3$) & exact & \Validated{} \\
Block-spark, broad peer (node 2, $21$ rows)& $1$ ($f_{\mathrm d}=0$) & exact & \Validated{} \\
Block-spark, narrow peer (one row)         & $2$ ($f_{\mathrm d}=1$) & exact & \Validated{} \\
Projector error (abs / rel)                & $1.447\times10^{-12}$ / $6.6\times10^{-10}$ & $<10^{-3}$ & \Validated{} \\
\bottomrule
\end{tabular}
\end{table}

\subsection{The independent oracle and non-circularity}
\label{sec:oracle}

The reference values are produced by \code{scripts/gen\_faultobs\_ref.py}, which reads the
observation \emph{rows} from the vendored \code{network.json} and recomputes each quantity
through a \emph{different} numerical stack: the parity projector from the SVD of
$\Wmat^{1/2}\Gmat$ (\code{numpy.linalg.svd}) rather than the engine's cyclic-Jacobi
eigensolver and spectral-sum pseudo-inverse; the detectability norm from a textbook Baarda
parity residual; the minimum detectable bias from the independent non-centrality
$\bm c^\top\Wmat\Pperp\bm c$; the rank and parity dimension from the numpy SVD singular
spectrum with a matched cutoff; and the block-spark from independent rank tests of the
projected peer blocks. The oracle's methods are genuinely independent of the crate's---SVD
versus cyclic Jacobi, a textbook parity norm versus the crate's $\bm c^\top\Wmat\Pperp\bm c$
evaluation, numpy SVD rank versus the crate's spectral rank---so agreement is a cross-check,
not a transcription.

The comparison is \emph{non-circular}: \code{reference.json} is derived from the input rows,
never from a Rust output. Byte-consistency is enforced by rounding every row float to six
decimals \emph{before} both writing \code{network.json} and reading it in the oracle, so both
implementations consume identical IEEE-754 doubles and the test isolates the two
\emph{algorithms}. The Rust integration test \code{tests/lunar\_faultobs\_reference.rs}
asserts agreement to relative and absolute error below $10^{-3}$ (integer counts exact) or
fails the build. The rank cutoff is chosen against the clean spectral gap of
Section~\ref{sec:unification}: the null directions sit at
$\lambda_{\text{null}}^{\max}=4.367\times10^{-15}$, the weakest observable direction at
$\lambda_{\text{obs}}^{\min}=5.455\times10^{-4}$, so a matched $10^{-9}$-relative cutoff
separates the two by eleven orders of magnitude and the counts are unambiguous.

\subsection{The honesty boundary: what is and is not validated}
\label{sec:boundary}

We state the scope with the discipline the paper is built on.
\begin{itemize}
  \item \textbf{Validated.} The linear algebra---the parity projector, the detectability
  norms, the minimum detectable bias, the rank / parity-dimension counts, and the Byzantine
  block-spark values---on real DE440-orientation rows, agreeing with an independent
  numpy/scipy oracle to below $10^{-3}$ (counts exact). This is a statement about a
  \emph{computation on these rows}, and nothing more.
  \item \textbf{Modelled, not validated.} The network geometry (eight nodes, the beacon and
  relay placement, the three integration windows, the weighting) is a representative
  construction. Every count read as an \emph{operational} tolerance---the Byzantine $f$ as a
  flight guarantee, the provider-count row, the protection-level crossing---is
  \Modelled{}; the validated object is the number the pipeline computes, not a mission
  assurance.
  \item \textbf{Proven, not externally validated.} The gauge dimension eight is the P4
  theorem \citep{baweja2026gauge}, imported; the undetectable characterization
  (Theorem~\ref{thm:undetectable}), the unification (Proposition~\ref{prop:unification}), the
  block-spark bound (Theorem~\ref{thm:byzantine}), the provider split
  (Proposition~\ref{prop:provider}), the holdover floor (Section~\ref{sec:holdover}), and the
  protection gap (Proposition~\ref{prop:protection}) are \Proven{} algebra, covered by module
  unit tests, and are \emph{not} claimed as externally validated.
\end{itemize}
This is the same firewall as the companion papers
\citep{baweja2026datum,baweja2026commonmode,baweja2026gauge}: exactly one externally
validated row, its scope stated precisely, everything else cited, proven, or modelled.

\section{The autonomous-holdover floor and the provider channel}
\label{sec:holdover}

The gauge \eqref{eq:gauge} has a temporal part, and that part is what a self-referential
network cannot hold: the common clock offset and, more consequentially, the common clock
\emph{rate}. This section states the holdover floor as a corollary of the temporal gauge,
frames it as the ensemble-time free parameter of the network, and connects the provider
channel of Section~\ref{sec:providersplit} to the external-tie requirement.

\subsection{The temporal gauge and the holdover floor}
\label{sec:holdovercor}

\begin{corollary}[Autonomous-holdover temporal-gauge floor]
\label{cor:holdover}
In the self-referential model, where elapsed time is measured by the network's \emph{own}
drifting clocks, the temporal part of the gauge $\Nsp(\Gmat)$ has dimension two---the common
offset $\dtau$ and the common rate $\dalpha$---so both are internally unobservable
(\Proven{} structure; the self-referential timekeeping setup is \Modelled{}). Adding a single
external rate tie reduces the temporal gauge to dimension one: the common rate becomes
observable and only the common offset survives (\Proven{}, by the PSD intersection law). No
internal measurement removes the common rate; per-node \emph{relative} rates are observable
across the integration windows, but the \emph{common} rate is not.
\end{corollary}

\begin{proof}
Adding a common offset $\dtau$ and a common rate $\dalpha$ to every node's clock changes no
inter-node differential observable, so $\bm e_{\dtau}^{\mathrm{common}}$ and
$\bm e_{\dalpha}^{\mathrm{common}}$ lie in $\Nsp(\Gmat)$ exactly (the temporal part of
\eqref{eq:gauge}), giving temporal-gauge dimension two; the engine reports
\code{rate\_in\_gauge}${}=\text{true}$, \code{temporal\_gauge\_dim}${}=2$. An external rate
tie contributes a PSD information block $\bm{I}_{\mathrm{rt}}$ that annihilates every vector
except those with a common-rate component; by the intersection law $\Nsp(\Gmat^\top\Wmat\Gmat
+\bm{I}_{\mathrm{rt}})=\Nsp(\Gmat^\top\Wmat\Gmat)\cap\Nsp(\bm{I}_{\mathrm{rt}})$, the common
rate is removed while the common offset survives, so \code{rate\_in\_gauge}${}=\text{false}$,
\code{temporal\_gauge\_dim}${}=1$. Multi-window observations constrain \emph{differences} of
node rates (per-node relative rates) but not their common part, which cancels differentially.
\end{proof}

The magnitude of this floor is the free-running physical drift of the network's clock
ensemble relative to Lunar Coordinate Time (TCL)---a \Modelled{} quantity set by the
oscillators, not by geometry. What is \Proven{} is the \emph{structure}: the common rate is a
genuine gauge direction with zero internal information (not merely a weakly-observed
direction), so an autonomous monitor cannot detect or correct it. This sharpens, in the
self-referential setting, the ensemble-time free-parameter idea that a timescale realized by
an ensemble of clocks has an irreducible common drift only an external reference fixes.

\subsection{The ensemble-time free parameter and Earth-link-loss holdover}
\label{sec:ensemble}

The holdover floor is precisely the contingency the paper is framed around. Under nominal
operation a lunar service is Earth-referenced---an Earth--Moon time comparison is exactly the
external rate tie of Corollary~\ref{cor:holdover}, pinning the common rate. When the Earth
link is lost, that tie is gone: the network keeps time with its own clocks, and the common
rate re-enters the gauge as an unobservable, uncorrectable drift. This is the temporal
counterpart of the datum-defect story---a lunar coordinate time realized internally has a
free rate parameter, just as a lunar frame realized internally has free rigid directions
\citep{baweja2026gauge}. Two current strands of lunar time standardization frame this floor,
and we keep them distinct. \emph{Lunar Coordinate Time} (TCL) is the IAU 2024 Resolution~II
\emph{coordinate} time---a TCG-analog defined for a lunar celestial reference system
\citep{iau2024ltc}---and is the ideal against which the network's internal drift is measured;
\emph{Coordinated Lunar Time} (LTC) is the \emph{realized} standard, a UTC-analog motivated by
the April 2024 US OSTP memorandum \citep{iau2024ltc,fienga2024moonlight}, that an operational
lunar service must actually disseminate. Either way the floor matters: a realized lunar
timescale must specify the external tie and cadence that bound the holdover, because no
internal monitor can. We state this as \emph{relevance}, not endorsement: nothing here is a
certified timescale realization or an agency-approved holdover budget.

\subsection{The provider channel}
\label{sec:providerchannel}

The provider fault is the structured, multi-node channel by which an autonomous network can
still cross-check its own frame and timescale without an Earth link---but only partly.
Proposition~\ref{prop:provider} splits a provider-wide bias into a common part in
$\rng(\Gmat)$ and a differential part outside it. The common part is exactly the class an
external tie is needed for: it is the shared component of the inter-ephemeris orientation and
position floor of P2 \citep{baweja2026interop}, and, being a whole-network datum move, it is
invisible to any number of mutually-consistent providers. The differential part is what
autonomy \emph{can} exploit: two providers detect their disagreement (but cannot say which is
wrong---the two-clock problem), and three or more identify the culprit through the block-spark
bound of Theorem~\ref{thm:byzantine}. The provider-count trend of Table~\ref{tab:provider}
($1\to0$, $2\to48$, $4\to72$ detectable cross-provider rows) is the \Modelled{} quantitative
face of this split. The design consequence---engineer at least three independent providers to
identify a faulty one, and retain an external tie for the common part---is
Section~\ref{sec:designlaws}.

\section{Design laws}
\label{sec:designlaws}

The proven structure of Sections~\ref{sec:obs}--\ref{sec:holdover} yields four design laws
for an autonomous lunar integrity architecture. Each law's \emph{logic} is \Proven{}; each
supporting \emph{magnitude} is \Modelled{}, drawn from the representative network and labelled
where it enters. None is a certified requirement, a technology-readiness claim, or an agency
endorsement; they are analysis-level guidance.

\begin{quote}
\textbf{Design law~1 (redundancy $\to$ Byzantine tolerance).} \emph{Cross-link and provider
redundancy raise the block-spark, and with it the number of Byzantine peers the network can
detect and identify.} By Theorem~\ref{thm:byzantine} the tolerances are $f_{\mathrm{detect}}=
\spark_{\mathrm b}-1$ and $f_{\mathrm{identify}}=\lfloor(\spark_{\mathrm b}-1)/2\rfloor$;
adding independent measurements raises $\spark_{\mathrm b}$ (six independent single-measurement
peers already give $\spark_{\mathrm b}=7$, $f_{\mathrm{detect}}=6$, $f_{\mathrm{identify}}=3$,
Table~\ref{tab:byzantine}). An autonomous architecture buys fault tolerance with engineered
redundancy, not with trust.
\end{quote}

\begin{quote}
\textbf{Design law~2 (limit single-node measurement dominance).} \emph{No single peer should
own a large share of a node's measurements.} A peer owning \emph{all} of a node's incident
rows has $\spark_{\mathrm b}=1$ and is undetectable (Section~\ref{sec:counts}): it can inject
a uniform self-bias equal to a shift of that node's own five local degrees of freedom, which
lies in $\rng(\Gmat)$ and produces zero residual. Broad measurement coverage by one peer is a
fragility, not a strength; the architecture must spread each node's measurements across
independent peers so that the union blocks stay full column rank.
\end{quote}

\begin{quote}
\textbf{Design law~3 (external tie for the common-rate and common-provider classes).}
\emph{The common clock rate and the common provider bias cannot be seen internally and
require an external time/frame tie.} The common rate is a temporal-gauge direction
(Corollary~\ref{cor:holdover}) and the common provider bias is in $\rng(\Gmat)$
(Proposition~\ref{prop:provider}); both have zero internal information. Autonomous RAIM cannot
substitute for an external tie for these classes; a standard must specify the tie and its
cadence (Section~\ref{sec:ensemble}).
\end{quote}

\begin{quote}
\textbf{Design law~4 (the protection gap: no finite alert limit is autonomous for
$\rng(\Gmat)$).} \emph{No finite alert limit protects the $\rng(\Gmat)$ fault class from an
internal monitor alone.} The Brown slope is infinite on $\rng(\Gmat)$
(Proposition~\ref{prop:protection}): those faults corrupt the state while leaving the parity
residual zero. A protection level built on the residual cannot bound them, so the architecture
must either exclude the $\rng(\Gmat)$ class by an external tie or accept it as an unbounded,
unmonitored risk.
\end{quote}

\paragraph{The protection gap, quantified (\Modelled{}).}
The finite side of design law~4 fixes the scale. For the representative network the worst
\emph{detectable} single-measurement fault has a finite Brown slope of $30.984$ (metres of
observable state error per unit weighted parity), while every in-range fault has slope
$\infty$ (Proposition~\ref{prop:protection}). Against a \emph{hypothetical} alert limit of
$50$~m---explicitly not an LNIS-certified value, used only to set a scale---the worst
detectable fault reaches the alert limit at a parity statistic of
$50/30.984\approx1.614$, so a residual monitor at that threshold protects the detectable
class. The in-range class, by contrast, reaches \emph{any} position error at parity zero
(silent): no finite alert limit protects it. This is the operational face of the one rank
deficiency---the same $\rng(\Gmat)$ that hides the fault (Theorem~\ref{thm:undetectable})
defeats the protection level (Proposition~\ref{prop:protection})---and the reason an
autonomous architecture must tie the frame and timescale externally for that class.

\paragraph{Representativeness ledger.}
Every magnitude above is \Modelled{} on the eight-node network of Section~\ref{sec:model}:
the slope $30.984$, the parity crossing $1.614$, the block-spark tolerances of
Table~\ref{tab:byzantine}, and the provider-count trend of Table~\ref{tab:provider} depend on
the representative geometry and would shift for a real constellation. What carries over
unchanged is the \Proven{} structure: the tolerance \emph{formulas}
($f=\spark_{\mathrm b}-1$, $\lfloor(\spark_{\mathrm b}-1)/2\rfloor$), the infinite slope on
$\rng(\Gmat)$, the undetectability of the common rate and common provider bias, and the
self-blindness of a broad-coverage peer. The alert limit is hypothetical; there is no
certified value, no flight datum, and no technology-readiness claim in any of the four laws.

\section{Discussion}
\label{sec:discussion}

\subsection{Honest novelty, restated against the prior art}
\label{sec:novelty}

We fix the boundary once more, now that the results are on the table, because the value of
this paper is its discipline about what is and is not new.

\emph{Versus rigidity-based autonomous fault monitoring \citep{iiyama2024,iiyama2025edm}.} The
recent rigidity and Euclidean-distance-matrix methods are the closest prior art: they monitor
an inter-satellite-linked constellation autonomously, and the EDM follow-on already detects
clock-jump faults and identifies a \emph{single} faulty satellite. Three distinctions remain
genuinely ours. (a)~Theirs is a simulation-validated \emph{algorithm}---a rigidity/GCEDM
singular-value test---whereas ours is an \emph{algebraic characterization} of which faults are
provably invisible (undetectable $\iff\rng(\Gmat)$) with provable bounds. (b)~They identify a
single culprit; Theorem~\ref{thm:byzantine} bounds worst-case \emph{colluding-coalition}
identification (block-spark $>2f$), a stronger adversary. (c)~They treat clock \emph{jumps} but
do not characterize the common offset and rate as an \emph{unobservable gauge with a holdover
floor} (Corollary~\ref{cor:holdover}), nor the coupled datum$\oplus$timescale unification
(Proposition~\ref{prop:unification}), nor the provider common/differential split. The deep
connection we concede: EDM/rigidity is intrinsically invariant to rigid transforms, so those
methods \emph{exploit} the spatial datum gauge as a tool (hence needing no precise ephemeris),
while we \emph{characterize} that invariance as the fault-hiding subspace and extend it to the
timescale. Our results and theirs are thus complementary.

\emph{Versus secure state estimation \citep{fawzi2014,shoukry2016,candes2005,donoho2003}.} The
sparse-observability / $2f$ bound is entirely theirs. Theorem~\ref{thm:byzantine} is its
geometric instantiation---the block-spark of the peer column-blocks of a physical lunar
$\Gmat$ carrying the coupled datum gauge. We claim the \emph{instantiation and its lunar
consequences} (the deficiency-of-five self-blindness, the provider split), not the bound.

\emph{Versus LiAISON \citep{hillborn2007}.} LiAISON obtains absolute state from crosslinks by
exploiting orbital dynamics over an arc; our analysis is the kinematic snapshot, whose gauge
is eight (Section~\ref{sec:scope}). The dynamical regime is the complementary way to reduce
that gauge, and we do not claim a dynamics-free absolute frame.

\emph{What is ours.} The coupled frame$\oplus$timescale gauge instantiated in fault
observability (the undetectable subspace is the measurement-space \emph{face} of the same rank
deficiency that produces the datum$\oplus$timescale gauge in parameter space); the geometric
secure-\emph{identification} bound for colluding peers;
the P1/P3/P4 gauge-unification; and the provider-fault taxonomy. Not ours: parity RAIM, the
MDB, the PL, the spark bound, the Byzantine model, and rigidity-based detection---all cited.

\subsection{Relation to the companion arc}
\label{sec:companions}

This is the fifth and closing paper of a program built on one open engine. P1 classifies the
spatial datum defect \citep{baweja2026datum}; P2 the cross-provider interoperability floor
\citep{baweja2026interop}; P3 the common-mode class an integrity monitor cannot see
\citep{baweja2026commonmode}; P4 the coupled frame$\oplus$timescale gauge
\citep{baweja2026gauge}. The present paper shows that these are not four separate gauges but
the fault-observability faces of one rank deficiency: the spatial datum (P1) and the temporal
gauge (P4) are the parameter-space $\Nsp(\Gmat)$; the common-mode class (P3) is a subset of
the measurement-space $\rng(\Gmat)$; and the interoperability floor (P2) is the common,
externally-tied part of the provider channel. Proposition~\ref{prop:unification} is the
statement that closes the loop. The five share the same firewall---exactly one externally
validated row per paper---and the same open substrate.

\subsection{Limitations}
\label{sec:limitations}

The firewall of Section~\ref{sec:anchor} is also the statement of limitations; we make them
explicit.
\begin{itemize}
  \item \textbf{Kinematic snapshot scope.} $\Gmat$ is the instantaneous geometry Jacobian;
  the gauge of eight is the snapshot datum defect. Orbital dynamics over an arc
  \citep{hillborn2007} can partially collapse $\Nsp(\Gmat)$ and is not modelled here. Our
  results are the geometry regime.
  \item \textbf{Representative geometry.} The eight-node network is a construction, not a
  flown constellation. The block-spark tolerances, the slope, the parity crossing, and the
  provider counts depend on it; the proven \emph{formulas} and structural results do not.
  \item \textbf{Counts are validated as computations, not as guarantees.} The Byzantine $f$,
  the provider-count row, and the protection crossing are outputs of the validated pipeline;
  read as operational flight tolerances they are \Modelled{}. There is no flight data.
  \item \textbf{First-order linearization.} The state, rows, and normal matrix are first
  order about a nominal; second-order geometric effects and non-Gaussian errors are outside
  the model.
  \item \textbf{No flight, certification, or TRL claim.} The alert limit is hypothetical, the
  holdover magnitude is a modelled oscillator drift, and nothing here is certified or endorsed
  by any agency.
\end{itemize}
None of these caveats touches the \Proven{} content---the projector survival, the
undetectable characterization, the unification, the block-spark bound, the holdover floor, and
the protection gap---or the single \Validated{} pipeline reproduction.

\subsection{Venue and standards relevance}
\label{sec:venue}

The results speak to autonomous integrity for lunar PNT (a natural fit for \emph{NAVIGATION}
or IEEE TAES, and for ION GNSS+ 2026) and to the current lunar reference-system and lunar-time
standardization---both the IAU~2024 coordinate time (Lunar Coordinate Time, TCL) and the
realized Coordinated Lunar Time (LTC) motivated by the 2024 US OSTP memorandum
\citep{iau2024ltc,fienga2024moonlight,lnis2023}. The
design laws are stated as relevance to those efforts---an autonomous architecture needs
engineered redundancy, distributed measurement coverage, an external tie for the common
classes, and an acknowledgement of the protection gap---not as certified requirements or as
any form of agency endorsement.

\section{Conclusion}
\label{sec:conclusion}

An autonomous lunar constellation, cut off from Earth, must certify its own integrity from
internal measurements alone, and this paper characterizes exactly what it can and cannot do.
The organizing result is a unification: the single rank deficiency of the self-referential
network Jacobian $\Gmat$ ($\rank\Gmat=32<40$) produces both the parameter-space
datum$\oplus$timescale gauge $\Nsp(\Gmat)$ of dimension eight (P1$\,\oplus\,$P4) and the
measurement-space blind fault subspace $\rng(\Gmat)$ of dimension $32$ (with the P3
common-mode class a subset), and an undetectable fault is precisely one absorbed into a state
error defined only modulo the gauge. We proved that a fault is undetectable if and only if it
lies in $\rng(\Gmat)$, with the classical parity characterization surviving the
rank-deficient, gauge-carrying $\Gmat$ through the pseudo-inverse projector.

On this footing we gave a geometric secure-identification bound: a coalition of $\le f$
colluding Byzantine peers is detectable when the parity-projected peer blocks are full column
rank (block-spark $>f$) and identifiable when block-spark $>2f$. This is the
sparse-observability / coding bound instantiated on the coupled-datum lunar geometry, not a
new bound, and it is explicitly distinct from---and not comparable to---the message-passing
consensus threshold; we cite the Byzantine model only. Its lunar consequences are ours: a
broad-coverage peer that owns all of a node's measurements is undetectable because it can
inject the node's own five-degree-of-freedom self-bias into $\rng(\Gmat)$, so broad coverage
is a fragility; and a provider-wide bias splits into an undetectable common part (needing an
external tie) and a detectable differential part that three or more providers can trace to a
culprit. We showed that the common clock rate is an irreducible temporal-gauge free
parameter---the autonomous-holdover floor, removed only by an external rate tie---and that no
finite alert limit protects the $\rng(\Gmat)$ class, whose Brown slope is infinite.

Exactly one claim is externally \Validated{}: the fault-observability linear-algebra
pipeline---parity projector, detectability norm, minimum detectable bias, rank and
parity-dimension counts, and Byzantine block-spark---reproduces an independent numpy/scipy
oracle to relative and absolute error below $10^{-3}$ on real DE440-orientation rows (defect
$8$, $\rank\Gmat=32$, parity dimension $52$, block-spark $7/1/2$; projector error
$1.447\times10^{-12}$). The geometry, every count read as an operational tolerance, and the
holdover magnitude are \Modelled{}; the rigidity, secure-estimation, and LiAISON prior art is
\Cited{} and bounded. Four design laws follow---engineer redundancy for tolerance, spread each
node's measurements, tie the common classes externally, and accept the protection gap---as
analysis-level guidance, with no flight datum, no certified alert limit, and no
technology-readiness claim. The proven and cited content carries over unchanged to any real
deployment; every number here is reproducible from the open engine by a single command.


\section*{Data and code availability}
\label{sec:availability}
The reference substrate for this paper is part of the open-source \code{kshana} engine
\citep{baweja2026kshana} (\url{https://github.com/ashfordeOU/kshana}; AGPL-3.0-only). The
self-referential per-node model \eqref{eq:state}, the differential observation-row builders
\code{differential\_range\_row} and \code{differential\_clock\_tie\_row}, the weighted parity
projector \code{parity\_projector} \eqref{eq:pperp}, the detectability test
\code{is\_detectable} ($\lVert\Pperp\bvec\rVert$), the minimum detectable bias \code{mdb}
(the $\bm c^\top\Wmat\Pperp\bm c$ form \eqref{eq:mdb}), the peer-signature and block-spark
routines \code{peer\_signature} and \code{byzantine\_bound}
(Definition~\ref{def:blockspark}, Theorem~\ref{thm:byzantine}), the temporal-gauge holdover
routine \code{holdover\_floor} (Corollary~\ref{cor:holdover}), and the Brown-slope /
protection-gap helper (Proposition~\ref{prop:protection}) are in the module
\code{lunar\_faultobs}. Every number in this paper is regenerated by
\begin{quote}
  \code{cargo run -{}-example p5\_autonomous\_fault\_observability}
\end{quote}
which prints Part~A (the detectability of the in-range and generic faults, the minimum
detectable bias, and the one-rank-deficiency counts $\rank\Gmat=32$, parity dimension $52$,
gauge dimension $8$ on the free-network $\Gmat$), Part~B (the Byzantine block-spark detect /
identify tolerances of Table~\ref{tab:byzantine} and the provider-count sweep of
Table~\ref{tab:provider}), and Part~C (the holdover temporal-gauge floor and the Brown-slope
/ protection-gap crossing of Section~\ref{sec:designlaws}), directly from the committed engine
and the vendored fixtures. No result is hard-coded in a print statement.

The single \Validated{} row---the fault-observability pipeline matching an independent
numpy/scipy oracle to relative and absolute error below $10^{-3}$ (integer counts exact) on
real DE440-orientation rows---is asserted in the continuous-integration suite by
\begin{quote}
  \code{cargo test -{}-test lunar\_faultobs\_reference}
\end{quote}
which fails the build on any mismatch. The fixtures are regenerated by
\code{cargo run -{}-example gen\_faultobs\_rows} (which builds the observation rows through the
real DE440 lunar physical-axis orientation path and rounds every float to six decimals)
followed by the numpy/scipy oracle \code{scripts/gen\_faultobs\_ref.py}; the vendored
artifacts are \code{tests/fixtures/faultobs/} (\code{network.json}, \code{reference.json},
\code{NOTICE.md}), whose notice records the exact scope of what is validated and what is
modelled. The DE440 lunar orientation is JPL's \citep{park2021de440}; the kernels are not
redistributed. A productized operational lunar integrity-monitoring capability---as opposed to
the open reference substrate that reproduces the paper's numbers---is a commercial capability
of \code{kshana-pro} and is not part of this open release.


\section*{Acknowledgements}
The numerical substrate underlying this work is part of the open-source
\code{kshana} engine. The lunar orientation data used here are the JPL DE440
development ephemeris. This research received no specific grant from any funding
agency in the public, commercial, or not-for-profit sectors.

\bibliographystyle{plainnat}
\bibliography{refs}

\appendix
\section{Methods}
\label{app:methods}

This appendix records the numerical methods behind the results: the parity-projector
construction through the pseudo-inverse, the $\Wmat$-self-adjointness note and the precise
minimum-detectable-bias form (the single- versus double-projector distinction), the
block-spark enumeration and its coalition cap, the gauge-generator construction, the
independence of the numpy/scipy oracle, and the rank tolerance. All are consistent with the
vendored provenance notice (Section~\ref{sec:availability}).

\subsection{The parity projector through the pseudo-inverse}
\label{app:projector}

The weighted parity projector \eqref{eq:pperp} is assembled as $\Pperp=\bm{I}-\Gmat
(\Gmat^\top\Wmat\Gmat)^{+}\Gmat^\top\Wmat$, where $(\Gmat^\top\Wmat\Gmat)^{+}$ is the
Moore--Penrose pseudo-inverse taken from the engine's information-matrix / spectral-sum
pseudo-inverse (\code{fim::crlb}), which remains finite even when $\Gmat^\top\Wmat\Gmat$ is
rank-deficient---as it is here, with the eight-dimensional gauge \eqref{eq:gauge}. Writing
$\bm M=\Wmat^{1/2}\Gmat$, the identity $\Wmat^{1/2}\Gmat(\Gmat^\top\Wmat\Gmat)^{+}\Gmat^\top
\Wmat^{1/2}=\bm M(\bm M^\top\bm M)^{+}\bm M^\top=\Pi_{\bm M}$ (the ordinary orthogonal
projector onto $\rng\bm M$) is what makes $\Pperp$ idempotent and $\Wmat$-self-adjoint through
the rank deficiency (Proposition~\ref{prop:projector}). The engine never forms $\Wmat^{1/2}$
explicitly; the oracle does, via the SVD of $\bm M$, which is the independence lever of
Section~\ref{app:oracle}.

\subsection{\texorpdfstring{$\Wmat$}{W}-self-adjointness and the minimum-detectable-bias form}
\label{app:mdb}

The projector $\Pperp$ is $\Wmat$-self-adjoint, $\Wmat\Pperp=\Pperp{}^\top\Wmat$, and
idempotent, $\Pperp{}^2=\Pperp$. These two facts fix the correct non-centrality of the Baarda
minimum detectable bias \eqref{eq:mdb} and dispose of a common error. The correct
non-centrality is the single-projector, weighted parity energy $\bm c^\top\Wmat\Pperp\bm c$.
It is \emph{equal} to the weighted energy of the projected column, for \emph{every} $\Wmat$:
\begin{equation}
  (\Pperp\bm c)^\top\Wmat(\Pperp\bm c)
  \;=\;\bm c^\top\Pperp{}^\top\Wmat\Pperp\bm c
  \;=\;\bm c^\top\Wmat\Pperp\Pperp\bm c
  \;=\;\bm c^\top\Wmat\Pperp\bm c,
  \label{eq:mdbid}
\end{equation}
using $\Pperp{}^\top\Wmat=\Wmat\Pperp$ (self-adjointness) then $\Pperp{}^2=\Pperp$
(idempotence). So $(\Pperp\bm c)^\top\Wmat(\Pperp\bm c)$ is \emph{not} a counterexample---it
is the same quantity. The genuinely \emph{wrong} form is the double-projector expression
\begin{equation}
  \bm c^\top\Pperp\Wmat\Pperp\bm c
  \;=\;\bm c^\top\Pperp\,(\Wmat\Pperp)\,\bm c
  \;=\;\bm c^\top\Pperp\Pperp{}^\top\Wmat\bm c
  \;=\;\bm c^\top(\Pperp\Pperp{}^\top)\Wmat\bm c,
  \label{eq:mdbwrong}
\end{equation}
which involves $\Pperp\Pperp{}^\top$ rather than $\Pperp$ and differs from $\bm c^\top\Wmat
\Pperp\bm c$ whenever $\Wmat\neq\bm{I}$ (because $\Pperp$ is $\Wmat$-self-adjoint, not
ordinary-symmetric, so $\Pperp\neq\Pperp{}^\top$ and $\Pperp\Pperp{}^\top\neq\Pperp$). Using
\eqref{eq:mdbwrong} in place of $\bm c^\top\Wmat\Pperp\bm c$ mis-scales the minimum detectable
bias under non-uniform weighting---the very regime of this paper, where $\sigma\in\{0.5,1,2\}$
gives $\Wmat\neq\bm{I}$. The engine's \code{mdb} uses the single-projector form
$\bm c^\top\Wmat\Pperp\bm c$; the oracle recomputes it independently, and the two agree, which
is part of what the \Validated{} row checks. The non-centrality $\lambda_0=17.075$ is the
standard one-degree-of-freedom value for $P_{\mathrm{fa}}=10^{-3}$, $P_{\mathrm{md}}=0.20$
\citep{baarda1968,teunissen2006}.

\subsection{Block-spark enumeration and the coalition cap}
\label{app:blockspark}

The block-spark \eqref{eq:blockspark} is computed by \code{byzantine\_bound} through exact
enumeration of peer-block unions: for increasing coalition size $k=1,2,\dots$, it forms every
union $\Bmat_T$ with $|T|=k$, projects it to $\Gbar_T=\Pperp\Bmat_T$, tests full column rank,
and returns the smallest $k$ at which a rank-deficient union appears. Enumeration is bounded
by a coalition cap \code{MAX\_COALITION} (the search stops once $\spark_{\mathrm b}$ is found
or the cap is reached), which is adequate because the network is small ($M=8$) and the
block-spark values of interest are small ($1$, $2$, $7$). The detect and identify tolerances
are read off as $f_{\mathrm{detect}}=\spark_{\mathrm b}-1$ and $f_{\mathrm{identify}}=
\lfloor(\spark_{\mathrm b}-1)/2\rfloor$ (Theorem~\ref{thm:byzantine}). Column-rank tests use
the same matched cutoff as the null-space rank (Section~\ref{app:rank}).

\subsection{The gauge-generator construction}
\label{app:gauge}

The eight analytic gauge generators of \eqref{eq:gauge} are built directly: three unit
translations (a common $\delta\bm p$ across all nodes), three infinitesimal rotations (the
skew-symmetric generators applied to the nodes' nominal positions), a common clock offset
(unit $\dtau$ across all nodes), and a common clock rate (unit $\dalpha$ across all nodes).
Applying $\Gmat$ to each generator yields a worst relative residual of $2.702\times10^{-17}$,
confirming that $\Nsp(\Gmat)$ contains the eight-dimensional datum$\oplus$timescale span
exactly (Section~\ref{sec:unification}); scale is deliberately \emph{not} a generator, because
under SI a differential range is length-bearing and scale is observable
\citep{baweja2026gauge}.

\subsection{Independence of the numpy/scipy oracle}
\label{app:oracle}

The oracle \code{scripts/gen\_faultobs\_ref.py} reproduces the pipeline through a
\emph{different} numerical stack: the parity projector from the SVD of $\bm M=\Wmat^{1/2}
\Gmat$ (\code{numpy.linalg.svd}) against the engine's cyclic-Jacobi eigensolver and
spectral-sum pseudo-inverse; the detectability norm from a textbook Baarda parity residual;
the minimum detectable bias from an independent evaluation of $\bm c^\top\Wmat\Pperp\bm c$;
the rank and parity dimension from the numpy SVD singular spectrum; and the block-spark from
independent rank tests of the projected peer blocks. The two backends are independent
implementations, so agreement to relative and absolute error below $10^{-3}$ (integer counts
exact) is a genuine cross-check, not a tautology. Non-circularity is structural: the oracle
reads the observation \emph{rows} from \code{network.json} and derives \code{reference.json}
from them, never from a Rust output; six-decimal rounding of the rows before both writing and
reading guarantees byte-identical inputs.

\subsection{The rank tolerance}
\label{app:rank}

Ranks (of $\Gmat$, of the projected peer blocks, and of the normal matrix) are taken with a
\emph{relative} tolerance $\text{rel\_tol}=10^{-9}$: a singular value or eigenvalue below
$\text{rel\_tol}\cdot\lambda_{\max}$ is treated as zero. This threshold is dimensionless and
invariant under positive column rescaling, so it is unaffected by the choice of measurement
units or preconditioning. For the free-network normal matrix the null directions sit at
$\lambda_{\text{null}}^{\max}=4.367\times10^{-15}$ while the weakest observable direction is
$\lambda_{\text{obs}}^{\min}=5.455\times10^{-4}$, a spectral gap of $1.249\times10^{11}$
against a largest eigenvalue $\lambda^{\max}=8.597\times10^{1}$; the matched cutoff separates
null from observable by eleven orders of magnitude, so the defect count of eight, the rank of
$32$, and the parity dimension of $52$ are unambiguous. The oracle applies the identical
matched cutoff to its SVD spectrum, so the integer counts agree exactly and the linear-algebra
comparison isolates the two algorithms, not two threshold choices.


\end{document}
